\subsection{Auswertung und Scoring}
Nach jedem Test hat die Automation die Testergebnisse ausgewertet und ein Scoring durchgeführt.
Mit diesem Scoring sollten aus den verschiedenen Messwerten leicht verständliche Kennzahlen erstellt werden, um die Umgebungen leichter vergleichen zu können. Dabei wurden für die Loadtests und Chaos-Tests verschiedene Scorings verwendet.

Das Scoring beginnt mit dem Wert 100 und wird reduziert, wenn die Messwerte eine geringere Zuverlässigkeit aufweisen.
Dabei hat jeder Messwert eine eigene Gewichtung, d. h., es wird unterschiedlich viel abgezogen. Zusätzlich hat jedes Ergebnis eines Messwert eine andere Gewichtung.
Ein Scoring kann dabei nie unter 0 gehen. Je höher das Scoring ist, desto besser ist die Zuverlässigkeit.

Beim Loadtest wurde für das Scoring die Error Rate als Maß für die Fehlertoleranz verwendet.
Zusätzlich wurden die Messwerte des CPU- und RAM-Verbrauchs des gesamten Systems sowie der Anwendung für die Wertung verwendet, um die Stabilität zu bewerten.
Auch die Anzahl der Neustarts floss in die Bewertung der Stabilität und der Fehlertoleranz ein.
In das Loadtest-Scoring floss auch die Performance minimal ein, um die Auswirkungen von der Zuverlässigkeit zu bewerten.

Beim Chaos-Test wird die Error Rate ebenso für die Fehlertoleranz beim Scoring verwendet.
Aus der Error Rate wird aber auch die Verfügbarkeit berechnet, damit diese im Scoring verwendet werden kann.
Weitere Metriken, die das Scoring beeinflussen, sind die Recovery Time und die Frage, ob die Umgebung sich wiederherstellen konnte.
Wenn das System es nicht geschafft hat, sich selbst wiederherzustellen, hat es direkt einen Score von 0 erhalten.